\section{Additional Results}\label{app:additional}

\subsection*{The compositional gap as a regression coefficient}

The shift-share reading of \cref{eq:shiftshare} can be restated as a regression, which shows that the differential survives within broad occupational families and not only across them. Stacking occupation-by-region observations,
\begin{equation}\label{eq:employment}
y_{k,r} = \alpha + \beta \Ebar_k + %\delta\, \Scot_r +
\zeta \bigl( \Ebar_k \times \Scot_r \bigr) + \phi_{s(k)} + u_{k,r},
\end{equation}
weighted by occupation-group employment, where $y_{k,r}$ is log occupational employment and, in a second specification, the employment share; $\phi_{s(k)}$ are one-digit SOC major-group fixed effects; and errors are clustered by occupation. Since $\Ebar_k$ does not vary by region, $\zeta$ is identified purely from the interaction of the common exposure score with the regional employment distribution, and a negative estimate means composition insulates Scotland. With major-group fixed effects, $\widehat{\zeta}$ is $-0.629$ (SE 0.241) for log employment and $-0.018$ (0.007) for the employment share. Within broad occupation groups, Scottish employment sits systematically further down the exposure gradient than rUK employment.

\subsection*{Measurement-error benchmark for the wage gradient}

Exposure enters the wage regression at the SOC minor-group level, so every worker in a minor group carries the same score and the measurement error is grouped rather than independent across observations. The classical correction $\widehat{\beta}^{\mathrm{EIV}}_1 = \widehat{\beta}^{\mathrm{OLS}}_1 / (1 - \widehat{\lambda})$, with reliability complement $\lambda = \sigma^2_u / \Var(E_j)$ and $\sigma^2_u = \widehat{\sigma}^2_{g(j)} / 10^4$, assumes independent homoskedastic error per observation, and the violation runs in the direction that overstates the correction. By the Frisch--Waugh--Lovell theorem the effective reliability is computed on the variation in $E$ left after residualising on the industry and year fixed effects and the worker controls of \cref{eq:wage}, so any control that explains part of $E$ worsens the attenuation. I report the classical figure as an upper benchmark on the attenuation rather than as the correction itself.

The formula is also a single-regressor result, which is why it cannot be read off term by term for the interaction. It applies once $\beta_3$ is read as the gap between the rUK exposure slope, $\beta_1$, and the Scottish one, $\beta_1 + \beta_3$. Each is a single mismeasured regressor within its own region, and since $\sigma^2_u$ is a property of the shared occupation score while $\Var(E)$ differs across regions only through employment composition, the two reliabilities are equal to a close approximation and the common factor $1/(1 - \widehat{\lambda})$ carries the interaction as well.

\begin{table}[h]
\centering
\caption{Scoring measurement error in the wage gradient: OLS and the classical benchmark}
\label{tab:eiv}
\small
\begin{tabular}{lccc}
\toprule
 & OLS & EIV benchmark & SE (OLS) \\
\midrule
Exposure ($\beta_E$) & 1.276 & 1.297 & 0.170 \\
Exposure $\times$ Scotland ($\beta_{E \times \Scot}$) &
  $-0.300$ & $-0.305$ & 0.065 \\
\bottomrule
\end{tabular}
\begin{minipage}{0.95\textwidth}
\vspace{0.4em}\footnotesize
\textit{Notes:} Worker-level regression of \cref{eq:wage} on the pooled 2022--2025 APS, $N = 176{,}444$, weighted by the APS income weight, with industry and year fixed effects and standard errors clustered by occupation, the level at which the exposure regressor varies; the occupation $\times$ region figure of 0.226 on the Scotland interaction is reported in \cref{sec:wages} as a comparison. The EIV benchmark applies $1/(1 - \widehat{\lambda})$ with $\widehat{\lambda} = 0.016$ to both the exposure slope and the interaction, the latter under the equal-reliability argument in the text.
\end{minipage}
\end{table}

The reliability complement is small, $\widehat{\lambda} = 0.016$, because the between-occupation variation in exposure dwarfs the scoring noise. The benchmark nudges the exposure slope up by 1.6\%, from 1.276 to 1.297, and the interaction by the same factor, from $-0.300$ to $-0.305$. Neither departure noted above is large enough to matter at that reliability. Overall, the correction is immaterial even under an estimator whose violated assumptions overstate it, so the grouped structure of the measurement error cannot rescue an attenuation this small.

%\subsection*{Why the coarser cluster gives the smaller standard error}

%Clustering by occupation gives a standard error of 0.065 on the Scotland interaction against 0.226 under occupation $\times$ region clustering, and the coarser level returning the \emph{smaller} figure is not an anomaly. In the cluster-robust sandwich the middle matrix accumulates $g_c = \sum_{i \in c} \tilde{x}_i \hat{u}_i$ over clusters, with $\tilde{x}$ the regressor residualised on the remaining covariates. For an occupation cluster $k$ the score splits as $g_k = g_{k,\Scot} + g_{k,\rUK}$, so
%\begin{equation}\label{eq:clustsplit}
%\E\bigl[ g_k g_k' \bigr]
%= \E\bigl[ g_{k,\Scot} g_{k,\Scot}' \bigr]
%+ \E\bigl[ g_{k,\rUK} g_{k,\rUK}' \bigr]
%+ 2\, \E\bigl[ g_{k,\Scot} g_{k,\rUK}' \bigr] ,
%\end{equation}
%whereas occupation $\times$ region clustering retains only the first two terms. For the interaction coefficient the residualised regressor loads with opposite signs on the two regions within an occupation, so a shock common to the occupation enters $g_{k,\Scot}$ and $g_{k,\rUK}$ with opposite signs and the cross term in \cref{eq:clustsplit} is negative. Occupation clusters nest occupation $\times$ region clusters and are the appropriate level; the finer clustering treats the identical score on the two sides of the border as two independent regressors and discards exactly the cross-region error covariance that the interaction contrast nets out.

